% ============================================================
% 中文幻灯片模板（slides-beamer）
% ------------------------------------------------------------
% 基于 Beamer + ctexbeamer，主题为模板自绘（无小众主题包依赖），
% 默认使用 TeX Live 自带的 fandol 字体。编译引擎：xelatex。
% 修改指南：
%   1. "演示信息"处填写标题、作者、单位、日期；
%   2. 每页用 \begin{frame}...\end{frame} 编写；
%   3. 讲到新章节时加 \section{...}，目录页会自动更新；
%   4. 主色在"主题设置"处修改（RGB 值），整份幻灯片同步变色；
%   5. 强调文字用 \alert{...}，要点块用 block 环境示例。
% ============================================================
\documentclass[
  UTF8,
  fontset=fandol,        % TeX Live 自带开源字体，跨平台可编译
  11pt,
  aspectratio=169,       % 16:9 宽屏（改 43 可回到 4:3）
]{ctexbeamer}

% ==================== 主题设置（自绘简洁主题） ====================
% 主色：改成你的机构色即可全局换色，例如深红 RGB(155,30,30)
\definecolor{maincolor}{RGB}{0, 82, 155}
\definecolor{lightbg}{RGB}{238, 243, 249}
\definecolor{graytext}{RGB}{90, 90, 90}

\usepackage{booktabs}   % 三线表
\usepackage{tikz}       % 简单示意图
\usetikzlibrary{positioning}  % 节点相对定位（below=of ... 语法）

\setbeamercolor{structure}{fg=maincolor}
\setbeamercolor{normal text}{fg=black!85}
\setbeamercolor{frametitle}{bg=maincolor, fg=white}
\setbeamercolor{title}{fg=maincolor}
\setbeamercolor{block title}{bg=maincolor, fg=white}
\setbeamercolor{block body}{bg=lightbg}
\setbeamercolor{alerted text}{fg=orange!85!black}
\setbeamercolor{footline}{bg=lightbg, fg=graytext}

% 页面标题条：左侧色块 + 简洁文字
\setbeamertemplate{frametitle}{
  \nointerlineskip
  \begin{beamercolorbox}[
    wd=\paperwidth, ht=1.6ex, dp=1.0ex, leftskip=1.0em
  ]{frametitle}
    \usebeamerfont{frametitle}\insertframetitle
  \end{beamercolorbox}
  \vspace{-0.3em}
}

% 页脚：短标题 + 页码
\setbeamertemplate{footline}{
  \begin{beamercolorbox}[
    wd=\paperwidth, ht=2.2ex, dp=1.0ex, leftskip=1.0em, rightskip=1.0em
  ]{footline}
    \usebeamerfont{footline}\insertshorttitle
    \hfill \insertframenumber\,/\,\inserttotalframenumber
  \end{beamercolorbox}
}

% 杂项：关闭导航符号，要点用圆点，公式用衬线字体
\setbeamertemplate{navigation symbols}{}
\setbeamertemplate{itemize items}[circle]
\setbeamertemplate{enumerate items}[default]
\setbeamerfont{frametitle}{size=\large, series=\bfseries}
\setbeamercovered{transparent}

% ==================== 演示信息 ====================
\title{示例研究汇报}
\subtitle{从问题到方法与实验}
\author{张~三}
\institute{示例大学~计算机学院}
\date{2026年8月31日}

\begin{document}

% ---- 封面页 ----
\begin{frame}[plain]
  \begin{center}
    \vspace{2.2em}
    {\usebeamerfont{title}\color{maincolor}\Huge\bfseries 示例研究汇报}\\[0.8em]
    {\color{maincolor}\rule{0.5\paperwidth}{1.2pt}}\\[0.8em]
    {\large 从问题到方法与实验}\\[2.5em]
    {\large 张~三}\\[0.4em]
    {\color{graytext} 示例大学~计算机学院 \quad 2026年8月31日}
  \end{center}
\end{frame}

% ---- 目录页：\section 会自动进目录 ----
\begin{frame}{汇报提纲}
  \tableofcontents
\end{frame}

% ==================== 第一部分 ====================
\section{研究背景}

\begin{frame}{研究背景与动机}
  % 每页要点建议不超过 4 条，一行一个观点
  \begin{itemize}
    \item 要点一：用一句话说明领域现状；
    \item 要点二：指出尚未解决的\alert{核心问题}；
    \item 要点三：说明解决该问题的价值与意义；
    \item 要点四：概括本文的切入点。
  \end{itemize}
  \begin{block}{一句话概括}
    把整个工作的核心思想浓缩成一句听众记得住的话。
  \end{block}
\end{frame}

% ==================== 第二部分 ====================
\section{方法设计}

\begin{frame}{总体框架}
  % columns 环境实现左右分栏
  \begin{columns}
    \begin{column}{0.52\textwidth}
      \begin{itemize}
        \item 模块一：数据采集与预处理；
        \item 模块二：核心模型（公式见下）；
        \item 模块三：评估与反馈。
      \end{itemize}
      \vspace{0.5em}
      核心目标函数：
      \[
        \mathcal{L}(\theta) =
        \frac{1}{N}\sum_{i=1}^{N} \ell\bigl(f(x_i;\theta),\, y_i\bigr)
      \]
    \end{column}
    \begin{column}{0.44\textwidth}
      % 用 TikZ 画一个简单流程示意（可整体删除换成图片 \includegraphics）
      \begin{center}
        \begin{tikzpicture}[
          node distance=0.55cm,
          box/.style={draw=maincolor, fill=lightbg, rounded corners=2pt,
                      minimum width=3.6cm, minimum height=0.75cm, font=\small},
          arr/.style={->, thick, maincolor},
        ]
          \node[box] (a) {输入数据};
          \node[box, below=of a] (b) {核心模型};
          \node[box, below=of b] (c) {结果输出};
          \draw[arr] (a) -- (b);
          \draw[arr] (b) -- (c);
        \end{tikzpicture}
      \end{center}
    \end{column}
  \end{columns}
\end{frame}

% ==================== 第三部分 ====================
\section{实验结果}

\begin{frame}{主要结果}
  % 三线表展示关键数据，突出自己方法的行
  \begin{table}
    \centering
    \begin{tabular}{lccc}
      \toprule
      方法 & 准确率（\%） & $F_1$（\%） & 参数量（M） \\
      \midrule
      基线方法 A & 85.2 & 84.1 & 21.3 \\
      基线方法 B & 87.6 & 86.9 & 34.8 \\
      \textbf{本文方法} & \textbf{90.4} & \textbf{89.7} & 28.5 \\
      \bottomrule
    \end{tabular}
  \end{table}
  \begin{itemize}
    \item 相比最优基线，准确率提升 \alert{2.8 个百分点}；
    \item 参数量更少，说明方法更轻量高效。
  \end{itemize}
\end{frame}

% ==================== 第四部分 ====================
\section{总结与展望}

\begin{frame}{总结}
  \begin{itemize}
    \item 回顾问题：一句重述要解决的问题；
    \item 回顾方法：一句话概括提出的方法；
    \item 回顾结果：一句话给出最有说服力的数字；
    \item 展望：说明下一步计划与开放问题。
  \end{itemize}
\end{frame}

% ---- 致谢 / 提问页 ----
\begin{frame}[plain]
  \begin{center}
    \vspace{3.5em}
    {\color{maincolor}\Huge\bfseries 谢谢聆听！}\\[1.2em]
    {\large \color{graytext} 欢迎批评指正 \quad Q \& A}
  \end{center}
\end{frame}

\end{document}
